%
% ********** Página de Rosto
%

% titlepage gera páginas sem numeração
\begin{titlepage}

\begin{center}

\small

% O comando @{} no ambiente tabular x é para criar um novo delimitador
% entre colunas que não a barra vertical | que é normalmente utilizada.
% O delimitador desejado vai entre as chaves. No exemplo, não há nada,
% de modo que o delimitador é vazio. Este recurso está sendo usado para
% eliminar o espaço que geralmente existe entre as colunas
\begin{tabularx}{\linewidth}{@{}l@{}C@{}r@{}}
% A figura foi colocada dentro de um parbox para que fique verticalmente
% centralizada em relação ao resto da linha
	\parbox[c]{3cm}{
		\includegraphics[width=\linewidth]{tex/00-estrutural/figuras/UFRN}
	} &
	
	\begin{center}
		\textsf{\textsc{
			Universidade Federal do Rio Grande do Norte\\
			Instituto do Cérebro\\
			Programa de Pós-Graduação em Neurociências
		}}
	\end{center} &
	
	\parbox[c]{3cm}{
		\includegraphics[width=\linewidth]{tex/00-estrutural/figuras/Ice}}
\end{tabularx}

% O vfill é um espaço vertical que assume a máxima dimensão possível
% Os vfill's desta página foram utilizados para que o texto ocupe
% toda a folha
\vfill

\LARGE

\textbf{Código de velocidade no hipocampo dorsal}

\vfill

\Large

\textbf{Me. José Henrique Targino Dias Góis}

\vfill

\normalsize

Orientador: Prof. Dr. Adriano Bretanha Lopes Tort
% Se não houver co-orientador, comente a próxima linha
%\\[2ex] Co-orientador: Prof. Dr. Beltrano Catandura do Amaral

\vfill

\hfill
\parbox{0.5\linewidth}{
	\textbf{Proposta de projeto de doutorado}
	apresentado ao Programa de Pós-Graduação em Neurociências da UFRN (área de concentração: neurocomputação, neuroengenharia e neuroterapia) como parte dos requisitos para obtenção do título de Doutor em Ciências.
}

\vfill

\large


Natal, RN, Abril de 2018

\end{center}

\end{titlepage}
